\chapter{Einleitung}
\label{ch:einleitung}

\section{Motivation und Szenario}
\label{sec:motivation}
In modernen Waschanlagen soll der Prozess der Programmauswahl und Abrechnung zunehmend automatisiert werden. Ein mögliches Szenario hierfür ist die Identifikation des gewählten Waschprogramms über einen \ac{QR}-Code, der vom Kunden auf einem Zettel hinter der Windschutzscheibe des Fahrzeugs platziert wird.

Die Herausforderung besteht darin, diesen \ac{QR}-Code zuverlässig zu erkennen, unabhängig von den widrigen Umgebungsbedingungen, die in einer Waschanlage herrschen. Dazu gehören nasse, seifige oder verschmutzte Scheiben, Spiegelungen sowie unterschiedliche Lichtverhältnisse. Zudem können sich weitere Objekte oder Zettel ohne \ac{QR}-Code hinter der Scheibe befinden, die nicht fälschlicherweise als relevanter Code erkannt werden dürfen. Die Erkennung soll dabei aus einer Entfernung von bis zu drei Metern erfolgen.

\section{Aufgabenstellung}
\label{sec:aufgabenstellung}
Das Ziel dieser Projektarbeit ist die Entwicklung eines robusten Systems zur optischen Erkennung und Analyse von \ac{QR}-Codes in dem beschriebenen Szenario. Die Aufgabe gliedert sich dabei in zwei Hauptbereiche:

\subsection{Teil 1: Erkennung mittels CNN}
Im ersten Schritt soll ein \ac{CNN} entwickelt werden, das entscheidet, ob sich ein \ac{QR}-Code im Bild befindet oder nicht. Hierbei sollen zwei Ansätze verfolgt und verglichen werden:
\begin{listNorm}
    \item Entwicklung und Optimierung einer eigenen \ac{CNN}-Architektur ohne die Verwendung vorgefertigter Modelle.
    \item Anwendung von Transfer Learning unter Verwendung von drei verschiedenen, bereits trainierten Modellen.
\end{listNorm}
Für den ersten Ansatz ist eine umfangreiche Hyperparameteroptimierung durchzuführen.

\subsection{Teil 2: Lokalisierung und Lesbarkeitsprüfung}
Im zweiten Teil der Aufgabe liegt der Fokus auf der Weiterverarbeitung der positiv klassifizierten Bilder. Es soll geprüft werden, ob der erkannte \ac{QR}-Code für einen handelsüblichen Scanner lesbar wäre.
\begin{itemize}
    \item \textbf{Visualisierung:} Lesbare \ac{QR}-Codes sollen im Bild mit einem grünen Rahmen, nicht lesbare mit einem roten Rahmen markiert werden.
    \item \textbf{Kameraausrichtung:} Basierend auf der Position und Verzerrung des \ac{QR}-Codes soll eine Empfehlung berechnet werden, um wie viel Grad eine darüber montierte Kamera geschwenkt werden müsste, um eine optimale Draufsicht (plan zum \ac{QR}-Code) zu erreichen.
\end{itemize}
